% Included by qstab_manual.tex. Public inputs only; no private classes.
\topic{Reproduce the scalar one-parameter article}{sec:oneparameterstudy}
One quaternion parameter has four real coordinates, but it is shared by
every coefficient of the perturbation polynomial. The order of multiplication
is part of the model:
\[
 p_L(s,t)=p_0(s)+t\,p_1(s),\qquad
 p_R(s,t)=p_0(s)+p_1(s)\,t,\qquad t\in\HH.
\]
Throughout this chapter, $p_0$ is Hurwitz and
$\deg p_1<\deg p_0$. The lower degree keeps the leading coefficient fixed,
so a continuous change of $t$ cannot send a zero through infinity.
These are continuous-time scalar-polynomial statements. The ordered
two-by-two transfer formula does not automatically give an exact radius
for a general uncertain quaternion state matrix.

\subsection{One example, selected examples, or the whole study}
The runner uses the source data of the scalar article, revision~7.
The keys are stable commands, not page numbers. From the public package root:
\begin{lstlisting}
qstab_setup('Examples',true);
run_qstab_oneparameter('list');
R = run_qstab_oneparameter('E1','Plots',true);
R.status
R.archiveFile
\end{lstlisting}
The default inputs use the native MATLAB \m{quaternion} class. To use only
ordinary component structures, add \m{'UseQuaternion',false}. The arithmetic
and mathematical input data are the same.

\begin{tabularx}{\linewidth}{@{}p{.15\linewidth}Y@{}}\toprule
Key & What the experiment distinguishes\\\midrule
E1 & Sensor versus actuator calibration; left and right quaternion gains.\\
E2 & A singular, nonzero perturbation frequency matrix.\\
E3, E4 & A sphere of nearest gains versus a poorly conditioned unique gain.\\
E5, E5V & A tangential first contact, with and without a common complex slice.\\
E5C & Multiplicity inflation when a real cubic is replaced by its norm.\\
E6, E7 & A nearly flat optimum, and large amplitude/frequency rescaling.\\
D11--D13 & Infinite ray margin, a later nonimaginary guardian candidate,
and a spherical contact.\\
A01 & Zero perturbation, nonmonic coefficients and nonunit directions.\\\bottomrule
\end{tabularx}

Select several keys with \m{'Cases',\{'E1','E3','E5V'\}}. Use
\m{'ExactCertificates',false} for the numerical study alone; no symbolic
toolbox is then needed. The default setting, \m{auto}, executes fresh exact
proofs when Symbolic Math Toolbox is usable and records their absence
otherwise. To verify the corrected symbolic path, first run its
focused tests, then require the complete article study:
\begin{lstlisting}
T = test_qstab_oneparameter('ExactCertificates',true);
T.archiveFile
assert(T.ok && T.complete, ...
    'Inspect the exact tests: %s',T.archiveFile)

R = run_qstab_oneparameter('all', ...
    'ExactCertificates',true,'Plots',true);
R.archiveFile
assert(R.ok && R.articleComplete, ...
    'Article checks are incomplete: %s',R.archiveFile)
\end{lstlisting}
These follow-up calls do not repeat the main public regression. They write
separate ZIPs, identified by \m{T.archiveFile} and \m{R.archiveFile}.
Retain both reports, including failed runs.

\m{R.ok} says that the checks actually executed passed.
\m{R.articleComplete} additionally requires every example and the required
fresh exact proofs. A numerical ray can correctly remain unresolved while
an independent exact calculation establishes its first contact. This
difference is retained in the report; a printed reference value is never
substituted for a computed answer.

\lead{Results and printed output}
The Command Window prints the selected keys, their status and the absolute
result paths. \m{R.outputDirectory} is a unique run directory below
\m{oneparameter\_results}. It contains \m{report.mat}, \m{report.txt},
\m{comparison.csv} and the requested figures. The runner writes a new
\m{qstab\_oneparameter\_results.zip} at the package root, replacing the
previous ZIP while retaining the unique run directories. Use
\m{R.archiveFile} to find the exact file to send. An error within an example
is recorded and does not prevent the final archive; invalid options or
unavailable explicitly requested input facilities fail during preflight.
\m{'Save',false,'Plots',false} gives an in-memory run without these files.

\subsection{Frequency, direction and global radius answer different questions}
At a fixed frequency $\omega$, \m{freqmargin} minimizes the gain norm over
the entire imaginary quaternion sphere of radius $\omega$. The global
radius then also minimizes over frequency:
\[
 \rho_S(\omega)=\min\{|t|:p_S(\cdot,t)
     \text{ has a zero class at }\pm\ii\omega\},
 \qquad r_S=\inf_{\omega\ge0}\rho_S(\omega).
\]
A ray calculation instead fixes a nonzero quaternion $u$ and varies only
the real number $\alpha$ in $t=\alpha u$. It returns the connected stable
interval containing zero; stable intervals can reappear farther away.

For E1 the data are $p_0=s^2+2s+1+\qj$ and $p_1=\unit s+\qk$.
The following public representation needs no polynomial class:
\begin{lstlisting}
P0 = struct('a',[1 2 1],'b',[0 0 0], ...
            'c',[1 0 0],'d',[0 0 0]);
P1 = struct('a',[0 0],'b',[0 1], ...
            'c',[0 0],'d',[1 0]);
F = qstab.family('oneparameter',P0,P1,'Side','right');
[rho,I] = qstab.freqmargin(F,sqrt(3)-1);
rho                  % approximately 0.378937381963012
I.parameterComponents

u = struct('a',0,'b',1,'c',0,'d',1);  % i+k, not a unit direction
[ab,J] = qstab.raymargin(F,u);
ab                   % [-Inf, 2-sqrt(3)]
[bounds,K] = qstab.stabilityradius(F);
bounds               % numerical lower/upper global radius bounds
K.status
\end{lstlisting}
The first two outputs have different meanings: \m{ab} contains signed
ray endpoints, whereas \m{bounds} encloses one nonnegative global radius.
For E1 the positive ray parameter is $2-\sqrt3$, while the corresponding
gain norm is $\sqrt2(2-\sqrt3)$ because $|u|=\sqrt2$.
If a family has finite parameter bounds, an unrestricted frequency or
global-radius query requires \m{'ExtendParameter',true}. A ray query keeps
the full-law stable interval and reports its intersection with the specified
parameter domain separately; a box edge is not a stability contact.

\subsection{Inspect the gain set and the evaluation zeros}
Different gains can have the same minimum norm. In E3,
$p_0=s^2+s+1$ and $p_1=1$, the minimum at $\omega=1/\sqrt2$ is attained on
\[
 t=-\tfrac12-v/\sqrt2,\qquad
 v\in\operatorname{Im}\HH,\quad |v|=1.
\]
Thus the answer is a two-sphere in four-dimensional gain space. It is not
the similarity class of an algebraic zero and should not be plotted as one.
\begin{lstlisting}
[rho,I] = qstab.freqmargin([1 1 1],1,1/sqrt(2));
G = I.geometry{1};
G.kind               % 'sphere'
G.centerComponents   % [-1/2;0;0;0]
G.sphereRadius       % 1/sqrt(2)
v = [0;1;0];
t = G.centerComponents + G.sphereRadius*G.basisComponents*v;
Z = I.zeroWitnesses{1};
Z.status
Z.rightEvaluationComponents
Z.leftEvaluationComponents
\end{lstlisting}
The three columns of \m{G.basisComponents} form an orthonormal basis of the
sphere's affine subspace; its rows are real, $\unit,\qj,\qk$ components.
The report returns one selected critical gain as well as this geometry.
\m{point} means a single reported nearest gain, \m{empty} means no attainable
contact at that frequency, and \m{undetermined} means the geometry was not
resolved reliably. A small singular-value gap does not establish exact
multiplicity. In particular, E4's small nonreal coefficient collapses E3's
sphere to a unique, sensitive gain. Inspect the gap, conditioning and
residual fields before comparing witness directions.

The two evaluation zeros may have different quaternion orientations while
representing the same algebraic class. They are reconstructed and checked
on the original critical polynomial. Use them only when \m{Z.status} is
\m{computed}. A \m{sphericalCandidate} reports a numerically vanishing
central quadratic remainder and requires an exact divisibility check for
a proof. The independent ordered complex-adjoint comparison in
\m{I.svdRadius} and \m{I.svdParameterComponents} can expose a sensitive gain
even when both radius calculations agree.

\subsection{Exact certificates are a separate, optional computation}
The ordinary numerical routines use floating point. In particular,
\m{K.status='converged'} means that its numerical bracket meets the requested
width tolerance; it does not certify every displayed digit.
With Symbolic Math Toolbox, proposed rational bounds can instead be proved
over the whole frequency half-line:
\begin{lstlisting}
[b,C] = qstab.radiuscertificate([1 1 1],1, ...
    {'0.8660','0.8661'});
assert(C.certified)
b
C.exactBounds
\end{lstlisting}
Decimal, fraction and exponent text denotes exact rational values:
\m{'0.8660'} is $433/500$, and \m{'1e-3'} is $1/1000$. The verifier checks
polynomial signs and exact real-root counts, including repeated roots and
flat extrema.
If a proposed interval fails, the returned status is a failed certificate,
not a claim of instability. This function verifies supplied bounds; it does
not silently replace them by the result of a numerical search.

For a tangency, a sign-change search can miss an isolated loss of strict
stability. E5 supplies a simple demonstration:

\begin{minipage}{\linewidth}
\begin{lstlisting}
P0 = struct('a',[1 2 1],'b',[0 0 0], ...
            'c',[1 0 0],'d',[0 0 0]);
P1 = struct('a',[-1 -1],'b',[0 0], ...
            'c',[-1 -1],'d',[0 0]);
[ab,N] = qstab.raymargin(P0,P1,1);  % may remain unresolved
[gammaBounds,E] = qstab.raycertificate(P0,P1,1);
assert(E.certified)
gammaBounds          % encloses the first positive contact gamma=1
E.lifting.status
\end{lstlisting}
\end{minipage}

Here \m{gammaBounds} encloses the first \emph{positive} ray endpoint.
It is not the signed interval \m{ab}; use direction $-u$ to certify the
opposite side. The exact method forms the norm-Hurwitz determinant and
isolates the first positive root of its square-free part. Even-multiplicity
contacts are retained. Verification of the original polynomial is reported
separately in \m{E.lifting}; the exact first-root proof is not confused with
a numerical evaluation residual.

\lead{What is exact input?}
A supplied double is interpreted as its exact stored binary rational.
Consequently the double \m{0.1} is not silently changed into $1/10$.
When the source specifies exact fractions or decimal text, supply
\m{'RationalData',D}, where \m{D.P0}, \m{D.P1}, and optionally
\m{D.Direction} are four-row rational component tables, with entries such
as \m{'1/10'} or \m{'0.1'} in cells. The explicit tables must round to
the numerical input. Symbolic inputs must already be exact rational numbers;
variable-precision values created by \m{vpa} are rejected. Increasing
\m{digits} cannot turn an approximate input into a rational proof.
Reports retain rational strings and ordinary MATLAB arrays, so they can
be saved without symbolic objects.

\subsection{Change frequency units without changing the gain}
E7 scales E1 by a large common amplitude and changes the frequency unit:
\[
 p_{j,K}(s)=10^{80}p_{j,E1}(s/10^6),\qquad
 \sigma=s/10^6.
\]
Both gain radii and critical gains are unchanged; physical contact
frequencies are multiplied by $10^6$. Dividing every coefficient by one
common amplitude alone does not fix the badly balanced powers of $s$.
The continuous numerical margin functions therefore accept
\m{'FrequencyScale','auto'} by default, or an explicit positive scale:
\begin{lstlisting}
K0 = struct('a',[1e80 2e74 1e68],'b',[0 0 0], ...
            'c',[1e80 0 0],'d',[0 0 0]);
K1 = struct('a',[0 0],'b',[0 1e74], ...
            'c',[0 0],'d',[1e80 0]);
[rho,I] = qstab.freqmargin(K0,K1,1e6*(sqrt(3)-1), ...
    'Side','right','FrequencyScale',1e6);
I.scaling
I.frequency          % original frequency units
I.workingFrequency   % divided by 1e6
\end{lstlisting}
Automatic scaling leaves ordinary well-balanced inputs alone. It records
every applied variable and common coefficient scale and preserves nonzero
coefficients. Ray directions are not normalized, and original frequency
units are restored in the returned diagnostics. These options apply to
continuous scalar polynomials, not to state-matrix or discrete-time calls.
Use the recorded scaling rather than disabling the nominal stability check
to get past an unresolved calculation.

\lead{Independent reproduction and provenance}
\m{reproducibility/qb01} contains source-linked rational inputs, historical
reference certificates and an explicitly invoked Python/SymPy replay.
MATLAB never launches Python automatically. A saved historical certificate
is a reference until a fresh verification is recorded. Likewise an auxiliary
Octave or Python execution is not a native MATLAB validation; the validation
ledger and each report identify the actual runtime and proof path.
